\section{Isolating a log canonical component}\label{sec:isolation}

The effective divisor constructed above may have several intersecting
components.  We now make one of them disjoint from the others, while
preserving the pullback of the nef divisor.  The construction adapts
\cite[Section~2.3]{Companion}: raise the boundary to a log canonical
threshold, then lower the coefficient of one component and run an MMP
that is trivial for the restored log canonical divisor.  The resulting
nefness forces the required disjointness.

\begin{proposition}\label{prop:isolation}
Let $k$ be an algebraically closed field of characteristic $p>3$.
Let $X$ be a terminal, $\mathbb Q$-factorial projective threefold, let
$(X,B)$ be log canonical with rational boundary, and let
\[
  0\ne N_0\sim m(K_X+B)
\]
be an effective nef Cartier divisor, where $m$ is a positive integer.
Assume that $N_0$ is numerically trivial on each component of its
support, and that every component of $B$ is either contained in or
disjoint from $\operatorname{Supp}N_0$.  Denote by $P_{\rm out}$ the union
of the latter components and put
\[
  U_X=X_{\rm reg}\setminus
       \bigl(\operatorname{Supp}N_0\cup P_{\rm out}\bigr).
\]
There exist a projective $\mathbb Q$-factorial klt threefold $Y$
birational to $X$, a rational boundary $G_Y$, an effective nef
$\mathbb Q$-Cartier divisor $N_Y$, and a prime divisor $S_Y$ with the
following properties.
\begin{enumerate}
\item The pair $(Y,G_Y)$ is log canonical, $S_Y$ has coefficient one
in $G_Y$, and $s:=\operatorname{mult}_{S_Y}N_Y$ is positive.  For
rational numbers $\tau>0$ and $\delta>0$,
\begin{equation}\label{eq:isolation-nef}
  K_Y+G_Y\sim_{\mathbb Q}\tau N_Y,
  \qquad K_Y+G_Y-\delta S_Y\ \text{is nef}.
\end{equation}
\item The divisor $N_Y$ is numerically trivial on every component of
its support.  Every component of $G_Y$ is either contained in or
disjoint from $\operatorname{Supp}N_Y$.
\item The divisor $S_Y$ is disjoint from every other component of
$N_Y$ and from $\operatorname{Supp}(G_Y-S_Y)$.  In particular,
\begin{equation}\label{eq:isolation-local}
  N_Y=sS_Y,\qquad G_Y=S_Y,\qquad
  K_Y+S_Y\sim_{\mathbb Q}\tau sS_Y
\end{equation}
on a neighborhood of $S_Y$, and $S_Y|_{S_Y}\equiv0$.
\item The birational correspondence identifies $U_X$ with an open
subset of $Y$.  On a smooth projective common resolution
$q:V_0\to X$, $r:V_0\to Y$ that is an isomorphism over this open,
\begin{equation}\label{eq:isolation-pullbacks}
  q^*N_0=r^*N_Y.
\end{equation}
The same equality holds after passing to any further common resolution.
\end{enumerate}
One may take $\tau=m^{-1}+\lambda$, where
$\lambda=\operatorname{lct}_{(X,B)}(N_0)\geq0$.
\end{proposition}

\begin{proof}
We first find a coefficient-one boundary component lying over
$\operatorname{Supp}N_0$.  On a log resolution of $(X,B+N_0)$, the log
canonical threshold is the minimum of the finitely many ratios
\[
  \frac{a(E,X,B)}{\operatorname{ord}_E(N_0)}
  \qquad\bigl(\operatorname{ord}_E(N_0)>0\bigr),
\]
where $a(E,X,B)$ denotes log discrepancy and the divisors considered
include the strict transforms of the components of $N_0$.
Thus $\lambda$ is finite and rational, possibly zero, and
$(X,B+\lambda N_0)$ is log canonical.  There is a log canonical place
$E$ of this pair for which $\operatorname{ord}_E(N_0)>0$.

Take a projective crepant $\mathbb Q$-factorial dlt modification
\[
  \mu:(W,G)\longrightarrow (X,B+\lambda N_0),
\]
with every extracted divisor having coefficient one; see
\cite[Theorem~2.27]{XuNuTwo}.  Set
\[
  N=\mu^*N_0,\qquad \tau=m^{-1}+\lambda.
\]
Then
\begin{equation}\label{eq:isolation-restored}
  K_W+G\sim_{\mathbb Q}\tau N.
\end{equation}
The divisor $N$ is nef and numerically trivial on every component of
its support: a curve in that support maps to a point or to a curve in
$\operatorname{Supp}N_0$, and the assertion follows by the projection
formula.

Every component of $G$ is either contained in or disjoint from
$\operatorname{Supp}N$.  For strict transforms of components of $B$
this follows from the hypothesis, while those coming from
$\lambda N_0$ are contained in the support.  An extracted divisor has center in
$\operatorname{Supp}N_0\cup P_{\rm out}$, since $X$ is klt and the
boundary vanishes outside this union.  Its irreducible center lies in
one of these two disjoint closed sets.  In the first case its order
on $N_0$ is positive, and in the second case the divisor is disjoint
from $\operatorname{Supp}\mu^*N_0$.

The center on $W$ of the place $E$ lies in $\operatorname{Supp}N$ and,
by crepancy and the dlt property, in $\lfloor G\rfloor$.  Choose a
component $S$ of $\lfloor G\rfloor$ containing this center.
The preceding separation then implies
$S\subset\operatorname{Supp}N$.  This argument also covers
$\lambda=0$; it does not require all components of $N$ to occur in
$G$.

The modification $\mu$ is an isomorphism outside
$\operatorname{Supp}N_0\cup P_{\rm out}$.  Indeed it has no exceptional
divisors there, and a projective small birational morphism to a
$\mathbb Q$-factorial variety is an isomorphism.  To recall the latter
fact, push down a relatively ample Cartier divisor.  A Cartier
multiple of its pushdown pulls back to the same multiple upstairs,
because there are no exceptional divisors.  It has degree zero on
every contracted curve, contradicting relative ampleness if such a
curve exists.

Choose a rational $\epsilon>0$ sufficiently small that
\begin{equation}\label{eq:isolation-effective}
  G-\epsilon S\geq0,\qquad \tau N-\epsilon S\geq0.
\end{equation}
Apply the partial MMP with lowered boundary $G-\epsilon S$ and
restoration divisor $\epsilon S$
\cite[Proposition~2.21, Definition~2.23, and Lemma~2.24]{XuNuTwo}.
Its steps are $(K_W+G-\epsilon S)$-negative and $(K_W+G)$-trivial.
The lowered pair is dlt, and its log canonical divisor is
pseudoeffective by \eqref{eq:isolation-restored} and
\eqref{eq:isolation-effective}.  The termination required here is
therefore supplied by \cite[Theorem~1.5]{XuNonvanishing}.
The threefold MMP package used in this application, including in
characteristic five, is recalled in \cite[Section~2.3]{Companion}.

The partial-MMP conclusion is either a Mori fiber contraction for the
lowered log divisor, trivial for the restored divisor, or nefness of the
restored divisor minus $\delta$ times the transform of $S$ for every
sufficiently small rational $\delta>0$. We first track $S$, the pullbacks
of $N$, and support separation through the steps, and then exclude the
fiber-type alternative. Write $W_i\dashrightarrow W_{i+1}$
for a step, with transforms $G_i,N_i,S_i$.  The equality
$K_{W_i}+G_i\sim_{\mathbb Q}\tau N_i$ persists by pushforward, so its
contracted ray $R_i$ satisfies
\begin{equation}\label{eq:isolation-ray}
  N_i\cdot R_i=0,\qquad S_i\cdot R_i>0.
\end{equation}
In particular $S_i$ cannot be the exceptional divisor of a divisorial
step: the negativity lemma would give a contracted curve of negative
$S_i$-degree.  Hence $S$ survives throughout the MMP.

Let $a:T\to W_i$ and $b:T\to W_{i+1}$ be a common resolution of a
step.  The divisor
\[
  F=a^*N_i-b^*N_{i+1}
\]
is $b$-exceptional.  If a curve on $T$ is contracted by $b$, its image
on $W_i$ is a point or a curve contracted to the common contraction
target.  Thus $F$ has degree zero on every such curve by
\eqref{eq:isolation-ray}.  The negativity lemma applied to $F$ and
$-F$ gives
\begin{equation}\label{eq:isolation-step-pullback}
  a^*N_i=b^*N_{i+1}.
\end{equation}
This applies to divisorial contractions and flips.  The same argument
for the full log divisors gives
\[
  a^*(K_{W_i}+G_i)=b^*(K_{W_{i+1}}+G_{i+1}).
\]
Consequently the restored pair remains log canonical.

Equation~\eqref{eq:isolation-step-pullback} also preserves nefness and
numerical triviality on the support.  For the latter assertion, lift
a curve in $\operatorname{Supp}N_{i+1}$ to a curve on $T$ dominating
it.  Its image on $W_i$ lies in $\operatorname{Supp}N_i$ or is a point:
the inverse image of the support of an effective $\mathbb Q$-Cartier
divisor is exactly the support of its pullback.  The projection
formula now gives degree zero.  Boundary components contained in
$\operatorname{Supp}N_i$ remain contained in the transformed support.
If a boundary component is disjoint from that support, its strict
transform on $T$ is disjoint from $a^*N_i$, hence from $b^*N_{i+1}$.
Since that strict transform maps onto its image on $W_{i+1}$, the
transformed boundary component is still disjoint from
$\operatorname{Supp}N_{i+1}$.

These steps also preserve the open complement of the support.
Indeed every contracted curve has negative degree against the
effective divisor $\tau N_i-\epsilon S_i$, so is contained in
$\operatorname{Supp}N_i$.  A connected positive-dimensional
projective fiber contains a curve through each of its closed points.
Connectedness of fibers therefore implies that a fiber meeting the complement of
$\operatorname{Supp}N_i$ consists of a single point and does not meet
the support.  Properness and normality show that the contraction is
an isomorphism from this complement onto an open subset of its
target.  For a flip, this open in the target is $\mathbb Q$-factorial,
being isomorphic to the corresponding open of $W_i$.  The new small
morphism is therefore also an isomorphism over it.

There is no Mori fiber output.  The effective representative
$\tau N_i-\epsilon S_i$ persists on each model, as does
$\tau N_i-tS_i$ for every rational $0<t\leq\epsilon$.
A curve in a fiber through a general point outside this effective
support has nonnegative degree, whereas the log divisor driving a
Mori fiber contraction has negative degree on every contracted
curve.  This rules out the fiber-type alternative for the lowered
class and for every such smaller perturbation of the restored class.

Let $Y$ be the resulting model.  The partial-MMP conclusion gives
\eqref{eq:isolation-nef} for some rational $0<\delta<\epsilon$.
Negative MMP
steps preserve the $\mathbb Q$-factorial dlt pair with lowered
boundary, so $(Y,G_Y-\epsilon S_Y)$ is dlt.  In particular its
$\mathbb Q$-factorial underlying variety $Y$ is klt: every place of
zero log discrepancy for a dlt pair has center in its reduced
boundary, and removing that effective $\mathbb Q$-Cartier boundary
strictly increases its discrepancy.  Here dlt
preservation is used for the lowered boundary; for the restored
boundary $G_Y$ we use the crepant comparison just proved.

It remains to see why nefness isolates $S_Y$.  If $T\ne S_Y$ is a
component of $N_Y$, numerical triviality of $N_Y|_T$ and
\eqref{eq:isolation-nef} imply that $-S_Y|_T$ is nef.  On the other
hand, $S_Y|_T$ is an effective $\mathbb Q$-Cartier divisor.  If it
were nonzero, its intersection with an ample Cartier divisor on the
integral projective surface $T$ would be positive, contradicting
anti-nefness.  Therefore $S_Y\cap T=\varnothing$.  Each other
component of $G_Y$ is either such a component of $N_Y$, or is
disjoint from its support.  This proves the asserted separation and
the local identities \eqref{eq:isolation-local}.  The equality
$N_Y|_{S_Y}\equiv0$ gives $S_Y|_{S_Y}\equiv0$.

Finally, all models share the stated open $U_X$.  Resolve the graph
of the correspondences among all the models, keeping this smooth
open unchanged, to obtain $V_0$.  Composing
\eqref{eq:isolation-step-pullback} with $N=\mu^*N_0$ proves
\eqref{eq:isolation-pullbacks}.  Pullback preserves this equality on
every further common resolution.
\end{proof}

For adjunction we replace the isolated log canonical pair by a dlt
pair.  The klt property of $Y$ ensures that all divisors extracted in
this last operation belong to the total transform of $S_Y$.

\begin{samepage}
\begin{corollary}\label{lem:dlt-boundary}
With the notation of Proposition~\ref{prop:isolation}, there is a
projective crepant $\mathbb Q$-factorial dlt modification
\[
  f:(Y_d,\Sigma)\longrightarrow (Y,S_Y)
\]
such that $\Sigma$ is reduced,
\begin{equation}\label{eq:dlt-support}
  \operatorname{Supp}\Sigma=\operatorname{Supp}f^*S_Y,
\end{equation}
and $f$ is an isomorphism off $S_Y$.  On a neighborhood of $\Sigma$,
\begin{equation}\label{eq:dlt-local}
  K_{Y_d}+\Sigma\sim_{\mathbb Q}\tau s f^*S_Y.
\end{equation}
Moreover, $f^*S_Y$ is numerically trivial on each component of
$\Sigma$, and $Y_d$ shares the open $U_X$ with $X$ and $Y$.
\end{corollary}
\end{samepage}

\begin{proof}
The pair $(Y,S_Y)$ is log canonical since $(Y,G_Y)$ is log canonical
and $G_Y\geq S_Y$.  Take a crepant dlt modification extracting only
log canonical places.  Its boundary $\Sigma$ consists of the strict
transform of $S_Y$ and its exceptional divisors, all with coefficient
one.  Every extracted divisor has center contained in $S_Y$:
outside $S_Y$ the pair is the klt variety $Y$.  Such a divisor has
positive order on $S_Y$, which proves \eqref{eq:dlt-support}.
There are no exceptional divisors off $S_Y$; since $Y$ is
$\mathbb Q$-factorial, the small-morphism argument in the preceding
proof shows that $f$ is an isomorphism there.

Crepancy and \eqref{eq:isolation-local} give
\eqref{eq:dlt-local}.  Every curve in $\Sigma$ maps to a point or to
a curve in $S_Y$, where $S_Y$ is numerically trivial.  The projection
formula proves the final numerical assertion.  Under the
identification in Proposition~\ref{prop:isolation}, $U_X$ is
disjoint from $S_Y$ and hence is unchanged by $f$.
\end{proof}
